\begin{table}[!htbp]\centering
\caption{Tres resultados centrales del contraste educativo}\label{tab:short_main}
\begin{threeparttable}\small
\setlength{\tabcolsep}{4pt}
\begin{tabular}{lrrrrr}
\toprule
Resultado & $N$ & DiD & IC 95\% & $p$ & $p$ previo \\
\midrule
Salario horario (log) & 69,534 & -0.021 & [-0.046, 0.005] & 0.106 & 0.001 \\
 &  & (0.012) &  &  &  \\
Formalidad (pp) & 181,317 & -4.642 & [-6.167, -3.117] & $<0.001$ & 0.148 \\
 &  & (0.739) &  &  &  \\
Autoempleo (pp) & 181,317 & 3.630 & [2.402, 4.859] & $<0.001$ & 0.003 \\
 &  & (0.595) &  &  &  \\
\bottomrule
\end{tabular}
\begin{tablenotes}\footnotesize
\item Coeficiente de baja educación por período posterior de la ecuación~\eqref{eq:short_did}. ENAHO 2021T1--2024T4, sin 2022T2. Salarios: asalariados privados con remuneración horaria válida. Formalidad y autoempleo: ocupados no públicos. Los indicadores binarios se expresan en puntos porcentuales (pp). Errores agrupados por departamento entre paréntesis; 25 grupos, ponderadores de encuesta. IC con distribución $t_{24}$. La última columna contrasta conjuntamente los coeficientes previos del estudio de eventos; no rechazar no demuestra tendencias paralelas.
\end{tablenotes}
\end{threeparttable}\end{table}
